% Paper 2 (EAAI) -- Results and Discussion. Impersonal. [R] = number from the
% evaluation-v1 runs (step 1: references and specialists, seed 42; seeds 43-46 for the
% dispersion). Design: docs/EVALUATION_PROTOCOL.md.
% Skeleton (user, 2026-11-17), every subsection on the four domains; no RQs, results by gap:
%   5.1 attack-based vs IID (FL = GL, Centralized as ceiling)       -- protocol 2.3, 3, 7
%   5.2 node loss: retention vs re-formed aggregate (static sweep)    -- protocol 5
%   5.3 mode switching on the Disaster-FD traces T1 and T2            -- protocol 9
%   5.4 window alarms at realistic prevalence                         -- protocol 6, 7
%   5.5 discussion, one paragraph per gap
% NO APPENDIX in this paper (user, 2026-11-17). k = 2 is a column of the 5.1 table.
% Precision at 1 % and 0.1 % prevalence opens 5.4. Pre-registered items still without a
% place (to decide): Local reference, per-type recall, Delta sensitivity (5.3), window
% sensitivities (5.4), dispersion over seeds 43-46.

\section{Results and Discussion}
\label{sec:results}

The results follow the four parts of the evaluation (Section~\ref{sec:evaldesign}), each in the four domains. Section~\ref{sec:res-partition} measures what the attack-based partition costs against an IID split of the same data, with a centralized model as a ceiling. Section~\ref{sec:res-nodeloss} removes nodes and compares the specialists retained by the survivors with an aggregate re-formed without the lost nodes. Section~\ref{sec:res-switching} replays the Disaster-FD trust levels and follows how each policy spreads the specialists of attacks that appear late. Section~\ref{sec:res-windows} turns the per-sample alerts into window alarms at realistic prevalence, and Section~\ref{sec:res-discussion} relates the findings to the gaps of Section~\ref{sec:intro}.
% [Check: the label of the introduction (sec:intro); "attacks that appear late" assumes the
%  Evaluation Design mentions the node blocks of 30 and 60 min, otherwise "spreads the
%  specialists over time".]

\subsection{Attack-based versus IID partitioning}
\label{sec:res-partition}

% [Text after the runs. Content to cover: same training rows, features and model budget
%  in the two partitions; FL = GL at full membership (retention); where the difference
%  between the partitions lies (recall or false positives) in each domain; Centralized as
%  a ceiling. Report every domain as measured, including any in which IID is not worse.
%  The k = 2 paragraph closes 5.1 (after the table). The rule itself stays in the
%  methodology.]

\begin{table}[htbp]
\centering
\caption{Detection under the IID and the attack-based partition, in the four domains
($N{=}14$ nodes, full test split of each domain, seed 42). Fusion by OR at 0.5, except in
the column $k{=}2$, which flags a sample when at least two specialists score at least 0.5,
on the same scores (not used by the system; it shows the veto of
Section~\ref{sec:fusion}).
\emph{IID}: the training rows split into 14 equal random parts, one binary XGBoost per node
with $\mathrm{round}(10S/14)$ trees. \emph{Attack-based}: one specialist per attack type,
each a binary XGBoost of 10 trees trained on its own attack rows against its node's slice
of benign rows. $S$ is the number of specialists (ERENO 14, CICIDS2017 17, Morris-4 35,
CICIoT2023 29), so both partitions hold about $10S$ trees on the same training rows and
features. With all nodes present, every node holds the same set of models in FL and in GL
(retention), so each partition yields one detector (FL = GL) in every column.
\emph{Centralized}: one
binary XGBoost of $10S$ trees on the pooled training rows of the 14 nodes; it requires the
data of all nodes in one place and is a ceiling, not a competitor. Precision and FPR at
the prevalence of each test split.}
\label{tab:partition-domains}
\small
\begin{tabular}{@{}llrrrr@{}}
\toprule
Domain & Metric (\%) & IID & Attack-based & Attack-based, $k{=}2$ & Centralized \\
\midrule
\multirow{4}{*}{ERENO}
 & F1-score  & [R] & [R] & [R] & [R] \\
 & Recall    & [R] & [R] & [R] & [R] \\
 & Precision & [R] & [R] & [R] & [R] \\
 & FPR       & [R] & [R] & [R] & [R] \\
\midrule
\multirow{4}{*}{CICIDS2017, improved}
 & F1-score  & [R] & [R] & [R] & [R] \\
 & Recall    & [R] & [R] & [R] & [R] \\
 & Precision & [R] & [R] & [R] & [R] \\
 & FPR       & [R] & [R] & [R] & [R] \\
\midrule
\multirow{4}{*}{Morris-4}
 & F1-score  & [R] & [R] & [R] & [R] \\
 & Recall    & [R] & [R] & [R] & [R] \\
 & Precision & [R] & [R] & [R] & [R] \\
 & FPR       & [R] & [R] & [R] & [R] \\
\midrule
\multirow{4}{*}{CICIoT2023}
 & F1-score  & [R] & [R] & [R] & [R] \\
 & Recall    & [R] & [R] & [R] & [R] \\
 & Precision & [R] & [R] & [R] & [R] \\
 & FPR       & [R] & [R] & [R] & [R] \\
\bottomrule
\end{tabular}
\end{table}
% Packages: booktabs, multirow.
% Dispersion over seeds 42-46: add "mean $\pm$ 95\% CI" in the cells or a footnote, once
% the runs show whether it is visible at two decimals.

Table~\ref{tab:partition-domains} also reports $k = 2$ on the same scores as OR, under the attack-based partition. In ERENO, where each attack type has two specialists, $k = 2$ changes the recall from [R] to [R]\,\% and the false-positive rate from [R] to [R]\,\%. In CICIDS2017, Morris-4 and CICIoT2023, where each type has a single specialist, it changes the recall from [R], [R] and [R]\,\% to [R], [R] and [R]\,\%, and the false-positive rate from [R], [R] and [R]\,\% to [R], [R] and [R]\,\%.
% [After the runs: one sentence on what this shows, as measured in each domain -- e.g. whether
%  the recall lost where a type has one specialist (veto, Section~\ref{sec:fusion}) outweighs
%  the false positives removed. Do not presuppose the direction.]

\subsection{Node loss: retention versus re-formed aggregate}
\label{sec:res-nodeloss}

% [Protocol section 5. Nodes removed in a fixed order (permutation, seed 42) from 14 to 3,
%  four domains. Retention (every surviving node keeps the specialists it received) against
%  the re-formed aggregate (only the specialists of the nodes still present). Overall and
%  per-type recall and FPR at every N; figure: recall vs N, one panel per domain, two
%  curves. The dynamic counterpart, driven by the failure detector, is T2 in 5.3.]

\subsection{Mode switching on the Disaster-FD traces}
\label{sec:res-switching}

% [Protocol section 9. T1 (run 1, dissemination) and T2 (scenario2, aggregator and region
%  loss; labelled synthetic if replayed from the expected trust levels). Blocks of nodes
%  8/12/14 at 0/30/60 min; Delta = 60 s (sensitivity 5-600 s in the appendix or a small
%  figure). Six policies: local, gossip fixed, federated fixed (retention / re-formed),
%  switching (retention = ReSIDS / re-formed). Content to cover: diffusion latency of the
%  late blocks per policy; recall on the late types per phase; reachable vs unreachable
%  nodes in T2; re-formed aggregate drops the blocks of nodes 2-7 at 75 min. Report as
%  measured that, in T2 between 60 and 75 min, gossip fixed reaches nodes 1 and 8 before
%  switching does (cutoff T_GL = 30 needs both ring neighbours). CICIoT2023: per-sample
%  metrics only (no windows), precision and F1 at 6.8 % derived.]

\subsection{Window alarms at realistic prevalence}
\label{sec:res-windows}

% [Protocol sections 6 and 7. Opens with per-sample precision at 1 % and 0.1 % prevalence
%  in the four domains (from R and FPR; 3/n bound when no false positive), which shows why
%  per-sample alarms are not usable in operation. Then windows in ERENO, CICIDS2017 and
%  Morris-4 (blocked-5 calibration): attack windows alarmed and false alarms, per type;
%  CICIDS2017 profile ablation as an extra row (count windows, no attribution). CICIoT2023
%  without windows, with the reason (rows shuffled; at 2.4 % benign almost every window
%  would hold an attack). Morris-4 and CICIDS2017 window sensitivities in the appendix.]

\subsection{Discussion}
\label{sec:res-discussion}

% [One paragraph per gap of the introduction, each drawing on 5.1-5.4 across the four
%  domains; the domain-profile gap is argued from what changed between domains (profile,
%  windows, calibration) and from the CICIDS2017 ablation, without a subsection of its own.
%  Closing: what held in every domain and what did not.]
